\documentclass[11pt]{article}

\usepackage[utf8]{inputenc}
\usepackage[T1]{fontenc}
\usepackage[margin=1in]{geometry}
\usepackage{amsmath,amssymb,amsthm,mathtools}
\usepackage{booktabs}
\usepackage{array}
\usepackage{enumitem}
\usepackage[hidelinks]{hyperref}

\theoremstyle{plain}
\newtheorem{theorem}{Theorem}[section]
\newtheorem{proposition}[theorem]{Proposition}
\newtheorem{lemma}[theorem]{Lemma}
\newtheorem{corollary}[theorem]{Corollary}
\theoremstyle{definition}
\newtheorem*{definition}{Definition}
\theoremstyle{remark}
\newtheorem*{remark}{Remark}

\DeclareMathOperator{\score}{score}
\DeclareMathOperator{\Bin}{Bin}
\DeclarePairedDelimiter{\abs}{\lvert}{\rvert}
\DeclarePairedDelimiter{\ceil}{\lceil}{\rceil}
\newcommand{\W}{\mathcal{W}}
\newcommand{\A}{\mathcal{A}}
\newcommand{\klo}{k_{\mathrm{lo}}}
\newcommand{\khi}{k_{\mathrm{hi}}}
\newcommand{\kaplo}{\kappa_{\mathrm{lo}}}
\newcommand{\kaphi}{\kappa_{\mathrm{hi}}}
\newcommand{\Sseg}{S_{\mathrm{seg}}}
\newcommand{\code}[1]{\texttt{\detokenize{#1}}}

\title{Proofs for the Alignment Kernel:\\
Diagonal Vote and Sliding-Window Smith--Waterman}
\author{TB Hardware Accelerator project, Algorithm focus}
\date{October 2026}

\begin{document}
\maketitle

\begin{abstract}
This note states and proves the claims behind the proposed architecture: a
diagonal vote followed by a sliding-window Smith--Waterman alignment. Each
result is labelled \emph{Proved} (a complete proof is given), \emph{Classical}
(a standard result, cited), \emph{Computed} (an exact value obtained from the
reference sequences by exhaustive calculation, with no sampling), or
\emph{Empirical} (supported only by simulated reads, and not a theorem).
Section~\ref{sec:open} lists what is not proved.
\end{abstract}

\section*{Summary of results}

\begin{center}
\begin{tabular}{@{}l>{\raggedright\arraybackslash}p{0.62\textwidth}l@{}}
\toprule
Result & Claim & Status \\
\midrule
Lemma~\ref{lem:paths} & $H(i,j)$ is the best score of a local alignment ending at $(i,j)$ & Classical \\
Prop.~\ref{prop:range} & $0 \le H(i,j) \le M\min(i,j)$; scores fit in 9 bits & Proved \\
Prop.~\ref{prop:diff} & Neighbouring cells differ by a bounded amount & Proved \\
Prop.~\ref{prop:cost} & Score equals $\tfrac{M}{2}(m+n)$ minus a shortest-path cost & Proved \\
Prop.~\ref{prop:work} & The window computes at most $mW$ cells & Proved \\
Prop.~\ref{prop:upper} & Window score never exceeds the full score & Proved \\
Prop.~\ref{prop:exact} & Window score equals the full score when an optimal path fits in the window & Proved \\
Theorem~\ref{thm:vote} & The vote lands on an acceptable diagonal when substitutions are below a margin & Proved \\
Cor.~\ref{cor:contain} & Then the window contains the true alignment & Proved \\
Tables~\ref{tab:margins}, \ref{tab:shares} & Values of the margin for the three loci & Computed \\
Cor.~\ref{cor:prob} & Probability bound for the vote under independent substitutions & Proved \\
Table~\ref{tab:hw} & Storage, cell and cycle counts & Proved \\
Prop.~\ref{prop:votecost} & The direct vote makes exactly $mn$ base comparisons & Proved \\
Prop.~\ref{prop:fft} & The vote can be computed in $O((m+n)\log(m+n))$ & Classical \\
Props.~\ref{prop:updates}, \ref{prop:compares} & Wavefront update and comparison bounds & Proved \\
Prop.~\ref{prop:segment} & Segment-vote band computes at most $mW + 2\Sseg G$ cells & Proved \\
Prop.~\ref{prop:gate} & Expected cost with a gate & Proved \\
\bottomrule
\end{tabular}
\end{center}

\section{Notation}

\begin{itemize}[nosep]
\item The read is $a = a_1 \dots a_m$ and the reference window is $b = b_1 \dots b_n$.
\item Scores: match $M > 0$, mismatch $x \le M$, and $g < 0$ per gapped base.
      The project uses $M = 2$, $x = -1$, $g = -2$, $m = 150$ and $n \le 180$.
\item $s(a_i, b_j) = M$ if $a_i = b_j$, and $s(a_i,b_j) = x$ otherwise.
\item Grid \emph{nodes} are pairs $(i,j)$ with $0 \le i \le m$ and $0 \le j \le n$.
      A node with $i, j \ge 1$ is also called a \emph{cell}.
\item The \emph{diagonal} of a node $(i,j)$ is $j - i$.
\item A \emph{path} is a sequence of steps between nodes. A diagonal step
      $(i-1,j-1) \to (i,j)$ scores $s(a_i,b_j)$. A down step $(i-1,j) \to (i,j)$
      and a right step $(i,j-1) \to (i,j)$ each score $g$. The score of a path is
      the sum of the scores of its steps, and the empty path scores $0$. A path is
      a local alignment of the substrings it spans.
\item $W = 2w + 1$ is the window width, $w$ its half-width, and $d_0$ the diagonal
      chosen by the vote.
\end{itemize}

\section{Smith--Waterman}

The recurrence below is the local alignment algorithm of Smith and
Waterman~\cite{smith1981} with a linear gap penalty.

\begin{definition}
$H(i,0) = H(0,j) = 0$, and for $i, j \ge 1$
\begin{equation}\label{eq:sw}
H(i,j) = \max\bigl\{\,0,\; H(i-1,j-1) + s(a_i,b_j),\; H(i-1,j) + g,\; H(i,j-1) + g \,\bigr\}.
\end{equation}
The alignment score is $S = \max_{i,j} H(i,j)$.
\end{definition}

\begin{lemma}[Classical]\label{lem:paths}
$H(i,j)$ is the maximum score over all paths that end at the node $(i,j)$, the
empty path included.
\end{lemma}

\begin{proof}[Proof sketch]
Induction on $i + j$. A non-empty path ending at $(i,j)$ ends with one of the
three steps, and the rest of it is a path ending at the corresponding
predecessor. Maximising over the last step and, by the induction hypothesis,
over the rest gives the three non-zero terms of \eqref{eq:sw}. The empty path
gives the term $0$. On the boundary, every non-empty path consists of gap steps
only and has a negative score.
\end{proof}

\begin{proposition}\label{prop:range}
$0 \le H(i,j) \le M \min(i,j)$. Consequently $0 \le S \le M\min(m,n)$, and an
unsigned score needs $\ceil{\log_2(M\min(m,n) + 1)}$ bits.
\end{proposition}

\begin{proof}
The lower bound is the term $0$ in \eqref{eq:sw}. For the upper bound use
Lemma~\ref{lem:paths}. A diagonal step increases both coordinates by one, so a
path ending at $(i,j)$ contains at most $\min(i,j)$ diagonal steps. Each scores
at most $M$, and every other step scores $g < 0$.
\end{proof}

With the project values, $S \le 300$ and $9$ bits suffice.

\section{Differences between neighbouring cells}

\begin{proposition}\label{prop:diff}
For all $1 \le i \le m$ and $1 \le j \le n$,
\begin{align}
g &\le H(i,j) - H(i,j-1) \le M - g, \label{eq:da}\\
g &\le H(i,j) - H(i-1,j) \le M - g, \label{eq:db}\\
x &\le H(i,j) - H(i-1,j-1) \le M. \label{eq:dc}
\end{align}
\end{proposition}

\begin{proof}
\emph{Lower bounds.} The recurrence \eqref{eq:sw} gives
$H(i,j) \ge H(i,j-1) + g$, $H(i,j) \ge H(i-1,j) + g$, and
$H(i,j) \ge H(i-1,j-1) + s(a_i,b_j) \ge H(i-1,j-1) + x$.

\emph{Upper bounds.} First extend \eqref{eq:da} and \eqref{eq:db} to the
boundary. For $i = 0$ and $j \ge 1$ we have $H(0,j) - H(0,j-1) = 0$, and for
$j = 0$ and $i \ge 1$ we have $H(i,0) - H(i-1,0) = 0$. Since $g \le 0 \le M - g$,
both inequalities hold there as well.

Now argue by induction on $i + j$. Fix $i, j \ge 1$ and assume the three upper
bounds hold at every node with a smaller coordinate sum. Consider which term
attains the maximum in \eqref{eq:sw}.

\emph{Case 0: $H(i,j) = 0$.} All three differences are at most $0$, because
$H \ge 0$ everywhere, and $0 \le M \le M - g$.

\emph{Case 1: $H(i,j) = H(i-1,j-1) + s(a_i,b_j)$.} Then
$H(i,j) \le H(i-1,j-1) + M$, which is \eqref{eq:dc}. The extended lower bound
\eqref{eq:db} at $(i,j-1)$ gives $H(i-1,j-1) \le H(i,j-1) - g$, hence
$H(i,j) \le H(i,j-1) + M - g$, which is \eqref{eq:da}. The extended lower bound
\eqref{eq:da} at $(i-1,j)$ gives $H(i-1,j-1) \le H(i-1,j) - g$, hence
$H(i,j) \le H(i-1,j) + M - g$, which is \eqref{eq:db}.

\emph{Case 2: $H(i,j) = H(i-1,j) + g$.} Then \eqref{eq:db} holds because
$g \le M - g$. By the induction hypothesis \eqref{eq:da} at $(i-1,j)$, or its
boundary form, $H(i-1,j) \le H(i-1,j-1) + M - g$, so
$H(i,j) \le H(i-1,j-1) + M$, which is \eqref{eq:dc}. Combining this with
$H(i-1,j-1) \le H(i,j-1) - g$ gives \eqref{eq:da}.

\emph{Case 3: $H(i,j) = H(i,j-1) + g$.} Symmetric to Case 2, with the roles of
$i$ and $j$ exchanged.
\end{proof}

\begin{corollary}\label{cor:bits}
A side-by-side difference \eqref{eq:da} or \eqref{eq:db} takes one of
$M - 2g + 1$ integer values and needs $\ceil{\log_2(M - 2g + 1)}$ bits. A
same-diagonal difference \eqref{eq:dc} takes one of $M - x + 1$ values and needs
$\ceil{\log_2(M - x + 1)}$ bits.
\end{corollary}

With the project values these are the ranges $-2,\dots,+4$ ($3$ bits) and
$-1,\dots,+2$ ($2$ bits). The simulated reads showed exactly these ranges.

\begin{remark}
Proposition~\ref{prop:diff} concerns the full matrix. In a window, a cell at the
edge has a neighbour outside the window that is read as $0$, and the bound does
not apply to that pair. The simulation measured only pairs with both cells
inside the window.
\end{remark}

\section{Score as a shortest-path cost}

For a path $P$ from $(i_0,j_0)$ to $(i_1,j_1)$ let $\mu$, $\nu$ and $\gamma$ be
its numbers of matches, mismatches and gap steps, and let
\[
u = i_0 + (m - i_1) + j_0 + (n - j_1)
\]
be the number of bases of $a$ and $b$ that $P$ leaves unaligned. Define the cost
\begin{equation}\label{eq:cost}
c(P) = (M - x)\,\nu + \Bigl(\frac{M}{2} - g\Bigr)\gamma + \frac{M}{2}\,u .
\end{equation}
All three coefficients are non-negative.

\begin{proposition}\label{prop:cost}
For every path $P$, $\score(P) = \frac{M}{2}(m+n) - c(P)$. Therefore
\[
S = \frac{M}{2}(m+n) - \min_{P} c(P).
\]
\end{proposition}

\begin{proof}
A diagonal step advances both coordinates and a gap step advances one of them,
so $(i_1 - i_0) + (j_1 - j_0) = 2(\mu + \nu) + \gamma$, and hence
$m + n = 2(\mu + \nu) + \gamma + u$. Substituting into \eqref{eq:cost},
\begin{align*}
\frac{M}{2}(m+n) - c(P)
 &= M\mu + M\nu + \frac{M}{2}\gamma + \frac{M}{2}u
    - (M - x)\nu - \frac{M}{2}\gamma + g\gamma - \frac{M}{2}u \\
 &= M\mu + x\nu + g\gamma ,
\end{align*}
which is the score of $P$. The cost is therefore a decreasing function of the
score over the same set of paths, the empty path included, so the maximum score
corresponds to the minimum cost.
\end{proof}

With the project values the costs are $0$ (match), $3$ (mismatch), $3$ (gapped
base) and $1$ (unaligned base), and $S = m + n - c_{\min}$. This is the quantity
a time-coded (``first arrival'') circuit would compute. The identity was also
checked numerically on $400$ random pairs.

\section{The sliding window}

\begin{definition}
The \emph{window} is the set of cells
\[
\W = \bigl\{(i,j) : 1 \le i \le m,\ 1 \le j \le n,\ \abs{j - i - d_0} \le w \bigr\}.
\]
$H_{\W}$ is defined by the recurrence \eqref{eq:sw} on the cells of $\W$ and is
$0$ at every other node. The window score is $S_{\W} = \max H_{\W}$.
\end{definition}

\begin{proposition}[Work]\label{prop:work}
$\abs{\W} \le mW$. The window therefore computes at most a fraction $W/n$ of the
$mn$ cells of the full matrix.
\end{proposition}

\begin{proof}
For each $i$, the condition $\abs{j - i - d_0} \le w$ allows at most
$2w + 1 = W$ values of $j$.
\end{proof}

For $m = 150$, $n = 180$ and $W = 17$ this is $2{,}550$ against $27{,}000$
cells, or $9.4\%$.

Call a path \emph{admissible} if every node on it, except possibly its start
node, is a cell of $\W$.

\begin{lemma}\label{lem:window}
For a cell $(i,j) \in \W$, $H_{\W}(i,j)$ is the maximum score over admissible
paths ending at $(i,j)$, the empty path included.
\end{lemma}

\begin{proof}
Induction on $i + j$. Let $p$ be a predecessor of $(i,j)$ and let $\sigma$ be
the score of the step from $p$ to $(i,j)$.

If $p \notin \W$, the recurrence uses $H_{\W}(p) = 0$ and the candidate value is
$\sigma$. That is the score of the one-step path from $p$, which is admissible
because $p$ is its start node. No longer admissible path passes through $p$,
because $p$ would then be a node other than the start node that lies outside
$\W$.

If $p \in \W$, the candidate value is $H_{\W}(p) + \sigma$. By the induction
hypothesis $H_{\W}(p)$ is the best score of an admissible path ending at $p$,
and appending the step keeps the path admissible. Conversely, an admissible path
whose last step comes from $p$ is an admissible path ending at $p$ followed by
that step.

Taking the maximum over the three predecessors and the empty path gives the
recurrence.
\end{proof}

\begin{proposition}\label{prop:upper}
$S_{\W} \le S$.
\end{proposition}

\begin{proof}
By Lemmas~\ref{lem:paths} and~\ref{lem:window}, $S_{\W}$ is a maximum over a
subset of the paths whose maximum is $S$.
\end{proof}

\begin{lemma}\label{lem:interval}
The diagonals of the nodes of a path form a set of consecutive integers
$\{\klo, \dots, \khi\}$, and $\khi - \klo \le \gamma$, the number of gap steps of
the path.
\end{lemma}

\begin{proof}
A diagonal step leaves the diagonal unchanged, a right step raises it by $1$ and
a down step lowers it by $1$. A sequence of integers that changes by at most $1$
at each step takes every value between its minimum and its maximum, and its
range is at most the number of steps that change it.
\end{proof}

\begin{proposition}[Exactness]\label{prop:exact}
Suppose $S > 0$ and some optimal path $P^{*}$ has diagonal range
$[\klo, \khi]$ with
\begin{equation}\label{eq:fits}
\khi - w \;\le\; d_0 \;\le\; \klo + w .
\end{equation}
Then $S_{\W} = S$. In particular this holds if $d_0$ is one of the diagonals of
$P^{*}$ and $P^{*}$ has at most $w$ gap steps.
\end{proposition}

\begin{proof}
An optimal path with positive score does not begin with a gap step, since
deleting that step would raise the score by $-g > 0$. So $P^{*}$ begins with a
diagonal step, and every node after its start node has $i \ge 1$ and $j \ge 1$;
that is, it is a cell. By \eqref{eq:fits} every diagonal $k$ of $P^{*}$
satisfies $\abs{k - d_0} \le w$, so those cells lie in $\W$ and $P^{*}$ is
admissible. Lemma~\ref{lem:window} gives $S_{\W} \ge \score(P^{*}) = S$, and
Proposition~\ref{prop:upper} gives equality.

For the last sentence, if $d_0 \in [\klo,\khi]$ and $\gamma \le w$, then
Lemma~\ref{lem:interval} gives $\khi - \klo \le w$, and therefore
$\khi - w \le \klo \le d_0 \le \khi \le \klo + w$.
\end{proof}

If $S = 0$ then $S_{\W} = 0$ as well, because $0 \le S_{\W} \le S$.

\begin{remark}
Proposition~\ref{prop:exact} is conditional. Whether the vote produces a $d_0$
that satisfies its hypothesis is the subject of Section~\ref{sec:vote}.
\end{remark}

\section{The diagonal vote}\label{sec:vote}

In this section indices start at $0$. The read is $r_0 \dots r_{m-1}$ and the
window is $b_0 \dots b_{n-1}$. Diagonal $k$ pairs $r_i$ with $b_{i+k}$, which is
the same numbering of diagonals as before.

\begin{definition}[Vote]
For $-(m-1) \le k \le n-1$ let
\begin{align*}
I(k) &= \{\, i : 0 \le i < m,\ 0 \le i + k < n \,\}, & N(k) &= \abs{I(k)},\\
V_r(k) &= \abs[\big]{\{\, i \in I(k) : r_i = b_{i+k} \,\}}, &
Z_r(k) &= (1+\lambda)\,V_r(k) - \lambda\,N(k),
\end{align*}
for a fixed penalty $\lambda \ge 0$. $N(k)$ is the number of compared positions
and depends only on $k$, $m$ and $n$. $Z_r(k)$ equals the number of matches
minus $\lambda$ times the number of mismatches. The vote returns a diagonal
$d_0$ that maximises $Z_r$. With $\lambda = 0$ it counts matches; with
$\lambda = 1$ it is matches minus mismatches.
\end{definition}

\begin{definition}[Read model]
Let $E$ be the reference genome around the window, with $b_j = E[\omega + j]$. A
\emph{template} $t_0 \dots t_{m-1}$ is given by a start $p$ and at most one
indel of length $g$ at read position $q$:
\begin{itemize}[nosep]
\item no indel: $t_i = E[p+i]$;
\item deletion: $t_i = E[p+i]$ for $i < q$, and $t_i = E[p+g+i]$ for $i \ge q$;
\item insertion: $t_i = E[p+i]$ for $i < q$, $t_i$ is unconstrained for
      $q \le i < q+g$, and $t_i = E[p-g+i]$ for $i \ge q+g$.
\end{itemize}
Let $F$ be the set of unconstrained positions, which is empty unless there is an
insertion. A read $r$ \emph{follows the template with at most $e$ substitutions}
if $r_i = t_i$ for all $i \notin F$ except at most $e$ positions.
\end{definition}

The \emph{true diagonals} are $k_1 = p - \omega$ for the part before the indel,
and $k_2 = k_1$ (no indel), $k_2 = k_1 + g$ (deletion) or $k_2 = k_1 - g$
(insertion) for the part after it. Write $\kaplo = \min(k_1,k_2)$ and
$\kaphi = \max(k_1,k_2)$. The \emph{true alignment} $T$ pairs each read base
with the reference base it was copied from, wherever that base lies in the
window. All its diagonals lie in $[\kaplo,\kaphi]$.

A diagonal $d$ is \emph{acceptable} if the window centred on it contains every
true diagonal:
\[
\A = \{\, d : \kaphi - w \le d \le \kaplo + w \,\}.
\]

For the template alone define the pessimistic and optimistic match counts
\[
V^{-}(k) = \abs[\big]{\{\, i \in I(k)\setminus F : t_i = b_{i+k} \,\}},
\qquad
V^{+}(k) = V^{-}(k) + \abs{I(k) \cap F},
\]
and $Z^{\pm}(k) = (1+\lambda)V^{\pm}(k) - \lambda N(k)$. The \emph{margin} of the
template is
\begin{equation}\label{eq:margin}
\Delta_\lambda = \max_{k \in \A} Z^{-}(k) \;-\; \max_{k \notin \A} Z^{+}(k).
\end{equation}
Here $k$ ranges over $-(m-1),\dots,n-1$. We assume $g \le 2w$, so that $\A$ is
not empty. If every diagonal is acceptable, the second maximum is over an empty
set and is taken to be $-\infty$.
The margin depends only on $E$, the window, $m$, $w$, $\lambda$ and the template
parameters $(p, q, g, \text{type})$. It can be computed before any read is seen.

\begin{lemma}\label{lem:perturb}
If $r$ follows the template with at most $e$ substitutions, then for every $k$
\[
V^{-}(k) - e \;\le\; V_r(k) \;\le\; V^{+}(k) + e .
\]
\end{lemma}

\begin{proof}
Split $I(k)$ into three sets. On positions outside $F$ where $r_i = t_i$, the
read and the template give the same indicator. On the at most $e$ substituted
positions, each indicator changes by at most $1$. On positions in $F$ the read
contributes between $0$ and $1$ each, which lies between the contribution
counted in $V^{-}$ (none) and in $V^{+}$ (all).
\end{proof}

\begin{theorem}[Vote]\label{thm:vote}
If $r$ follows the template with at most $e$ substitutions and
\begin{equation}\label{eq:votecond}
2\,(1+\lambda)\,e \;<\; \Delta_\lambda ,
\end{equation}
then every maximiser of $Z_r$ is acceptable.
\end{theorem}

\begin{proof}
Let $k_A \in \A$ attain $\max_{k \in \A} Z^{-}(k)$. By Lemma~\ref{lem:perturb},
$Z_r(k_A) \ge Z^{-}(k_A) - (1+\lambda)e$. For any $k \notin \A$,
Lemma~\ref{lem:perturb}, the definition \eqref{eq:margin} and the hypothesis
\eqref{eq:votecond} give
\[
Z_r(k) \le Z^{+}(k) + (1+\lambda)e
       \le Z^{-}(k_A) - \Delta_\lambda + (1+\lambda)e
       <   Z^{-}(k_A) - (1+\lambda)e
       \le Z_r(k_A).
\]
So no diagonal outside $\A$ maximises $Z_r$.
\end{proof}

\begin{corollary}[The window contains the true alignment]\label{cor:contain}
Under the hypothesis of Theorem~\ref{thm:vote}, the window centred on $d_0$
contains every diagonal of the true alignment $T$, and
\[
\score(T) \;\le\; S_{\W} \;\le\; S .
\]
\end{corollary}

\begin{proof}
$d_0 \in \A$ means $d_0 - w \le \kaplo$ and $\kaphi \le d_0 + w$, so every
diagonal of $T$ is within $w$ of $d_0$. Take $T$ restricted to the grid,
starting at its first paired base. Its nodes after the start node are cells on
those diagonals, so the path is admissible, and Lemma~\ref{lem:window} gives
$\score(T) \le S_{\W}$. The second inequality is Proposition~\ref{prop:upper}.
\end{proof}

\begin{remark}
This is the guarantee that matters for calling mutations: the window cannot lose
the alignment the read actually came from. It does not assert $S_{\W} = S$; see
Section~\ref{sec:open}.
\end{remark}

\subsection{Margins for the three loci (Computed)}\label{sec:margins}

The script \code{reference/vote_margin.py} evaluates $\Delta_\lambda$ exactly
for $m = 150$, $w = 8$, every read start for which the read covers the hotspot
centre with at least $20$~bp on each side ($111$ starts per locus), every indel
position, every length $g = 1,\dots,8$ and both indel types. That is $260{,}739$
templates per locus and $782{,}217$ in total. The windows are $180$~bp for rpoB
and $120$~bp for katG and inhA, taken from NC\_000962.3. The window coordinates
are provisional.

Table~\ref{tab:margins} gives the minimum margin and the largest $e$ that
satisfies \eqref{eq:votecond}. Table~\ref{tab:shares} gives the share of
templates for which the vote is guaranteed with at least $e$ substitutions.
These shares weight every template equally; they are not probabilities over real
reads.

\begin{table}[ht]
\centering
\caption{Minimum margin $\Delta_\lambda$ over all templates, and the largest
tolerated number of substitutions $e$.}\label{tab:margins}
\begin{tabular}{@{}llrrrr@{}}
\toprule
 & & \multicolumn{2}{c}{$\lambda = 0$} & \multicolumn{2}{c}{$\lambda = 1$} \\
\cmidrule(lr){3-4}\cmidrule(l){5-6}
Locus & Case & $\min\Delta$ & $e$ & $\min\Delta$ & $e$ \\
\midrule
rpoB & no indel            & 56   & 27   & 106 & 26 \\
rpoB & deletion $\le 8$    & 5    & 2    & 12  & 2  \\
rpoB & insertion $\le 8$   & $-4$ & none & 10  & 2  \\
katG & no indel            & 32   & 15   & 78  & 19 \\
katG & deletion $\le 8$    & 1    & 0    & 9   & 2  \\
katG & insertion $\le 8$   & $-9$ & none & 4   & 0  \\
inhA & no indel            & 36   & 17   & 78  & 19 \\
inhA & deletion $\le 8$    & 5    & 2    & 14  & 3  \\
inhA & insertion $\le 8$   & $-4$ & none & 9   & 2  \\
\bottomrule
\end{tabular}
\end{table}

\begin{table}[ht]
\centering
\caption{Share of all $782{,}217$ templates for which the vote is guaranteed
with at least $e$ substitutions.}\label{tab:shares}
\begin{tabular}{@{}lrrrrrr@{}}
\toprule
$e$ & 0 & 1 & 2 & 3 & 5 & 8 \\
\midrule
$\lambda = 0$ & 99.91\% & 99.83\% & 99.70\% & 99.50\% & 98.79\% & 96.55\% \\
$\lambda = 1$ & 100\%   & $>99.99\%$ & $>99.99\%$ & 99.95\% & 98.99\% & 92.61\% \\
\bottomrule
\end{tabular}
\end{table}

The tables support four statements.
\begin{enumerate}
\item \emph{With the plain match count ($\lambda = 0$) the vote is not
      guaranteed.} For some insertion templates the margin is negative, so a
      wrong diagonal can win even with no substitution at all. The cause is that
      a diagonal on which more of the read overlaps the window accumulates
      chance matches over more positions. The sequence is GC-rich, and about
      $40\%$ of unrelated positions match.
\item \emph{With $\lambda = 1$ every template has a positive margin.} A read
      with no substitutions always votes for an acceptable diagonal, for every
      start and every single indel of up to $8$~bp at all three loci.
\item \emph{Reads without an indel are safe by a wide margin} under either
      statistic: at least $15$ substitutions are tolerated.
\item \emph{Indel reads have thin worst-case margins.} A few templates tolerate
      no substitution, or only two. The bound is adversarial: it assumes that
      every substitution removes a match from the true diagonal and adds one to
      the same wrong diagonal. It is therefore conservative.
\end{enumerate}

The hardware cost of $\lambda = 1$ is small. $Z(k) = 2V(k) - N(k)$, and $N(k)$
is a function of the shift only, so it needs one counter and one subtraction.

\begin{corollary}[Independent substitutions]\label{cor:prob}
Suppose each of the $m$ read bases is substituted independently with probability
$\varepsilon$, and the template tolerates $e^{*}$ substitutions. Then
\[
\Pr[\text{the vote is not acceptable}]
\;\le\; \Pr\bigl[\Bin(m,\varepsilon) > e^{*}\bigr]
= \sum_{i > e^{*}} \binom{m}{i}\,\varepsilon^{i}(1-\varepsilon)^{m-i}.
\]
\end{corollary}

\begin{proof}
By Theorem~\ref{thm:vote}, the vote can fail only if more than $e^{*}$
substitutions occur.
\end{proof}

For $m = 150$ the bound evaluates as follows.
\begin{center}
\begin{tabular}{@{}rlccc@{}}
\toprule
$e^{*}$ & Applies to & $\varepsilon = 0.005$ & $\varepsilon = 0.01$ & $\varepsilon = 0.02$ \\
\midrule
19 & every read without an indel, $\lambda = 1$
   & $1.9\times10^{-22}$ & $1.0\times10^{-16}$ & $3.1\times10^{-11}$ \\
3  & $99.95\%$ of templates, $\lambda = 1$
   & $7.1\times10^{-3}$ & $6.5\times10^{-2}$ & $3.5\times10^{-1}$ \\
\bottomrule
\end{tabular}
\end{center}
The second row is weak because the bound is adversarial, not because the vote
fails that often. Here $\varepsilon$ must include true SNPs as well as
sequencing errors.

\paragraph{Empirical check.}
On $16{,}500$ simulated reads (SNPs, $0.5\%$ errors, indels of up to $8$~bp),
the $W = 17$ window disagreed with full Smith--Waterman on $1$ read with either
statistic, although the two statistics chose different diagonals on $143$ reads.
The change to $\lambda = 1$ therefore buys a guarantee; it showed no measurable
difference in simulation.

\section{Storage, cells and cycles}

Table~\ref{tab:hw} follows by counting.

\begin{table}[ht]
\centering
\caption{Hardware quantities for $m = 150$, $n = 180$, $W = 17$.}\label{tab:hw}
\begin{tabular}{@{}lccc@{}}
\toprule
Quantity & Full SW, systolic & Window SW & Diagonal vote, streaming \\
\midrule
DP cells in hardware & $\min(m,n) = 150$ & $W = 17$ & none ($m$ two-bit comparators) \\
Cycles per read      & $m + n - 1 = 329$ & $m = 150$ & $m + n - 1 = 329$ \\
Traceback bits       & $2mn = 54{,}000$  & $2mW = 5{,}100$ & none \\
\bottomrule
\end{tabular}
\end{table}

\begin{itemize}
\item Without pipelining, one read takes $(m + n - 1) + m = 479$ cycles through
      the vote and the window. That is more than the $329$ cycles of a full
      systolic array.
\item With the two stages pipelined, throughput is set by the vote: one read per
      $m + n - 1$ cycles, the same as the full array.
\item For the dynamic-programming part alone, the product of cells and cycles
      falls from about $\min(m,n)\,(m+n)$ to $mW$.
\end{itemize}
The design therefore reduces area and storage. It does not reduce the number of
clock cycles per read.

\section{Cost of the vote}

\begin{proposition}\label{prop:votecost}
Computing $V(k)$ for every $k$ by direct comparison uses exactly $mn$ base
comparisons.
\end{proposition}

\begin{proof}
$\sum_k N(k)$ counts the pairs $(i,j)$ with $0 \le i < m$ and $0 \le j < n$,
each exactly once, on its diagonal $k = j - i$.
\end{proof}

In a sequential model the direct vote is therefore $\Theta(mn)$, the same order
as full Smith--Waterman. Its advantage there is the constant: a two-bit
comparison against a cell update. With machine words of $\omega$ bits the
comparisons pack into $O(mn/\omega)$ word operations.

\begin{proposition}[Classical]\label{prop:fft}
All values $V(k)$ can be computed in $O\bigl((m+n)\log(m+n)\bigr)$ time.
\end{proposition}

\begin{proof}[Proof sketch]
\[
V(k) = \sum_{c \in \{\mathrm{A},\mathrm{C},\mathrm{G},\mathrm{T}\}}
       \sum_{i} [r_i = c]\,[b_{i+k} = c].
\]
For each letter $c$ the inner sum is a cross-correlation of two $0/1$ vectors,
which a fast Fourier transform computes for all $k$ at once. Four transforms
suffice. This is the match-counting technique of Fischer and
Paterson~\cite{fischer1974}.
\end{proof}

Hence the vote followed by the window costs $O\bigl((m+n)\log(m+n) + mW\bigr)$
in a sequential model, which is below $\Theta(mn)$ when $W = o(n)$. The hardware
does not use this method, so this bound is not a claim about the hardware.

\section{The wavefront alternative}

This section concerns Myers's algorithm~\cite{myers1986} for two strings of
lengths $N_a$ and $N_b$. Let $D$ be the minimum number of single-base insertions
and deletions that turn one string into the other; a substitution counts as $2$.
Place the strings on a grid with coordinates $(x,y)$ and number the diagonals by
$k = x - y$. A horizontal step raises $x$ and $k$, a vertical step raises $y$
and lowers $k$, and a diagonal step along a match is free.

The algorithm works in rounds $d = 0, 1, 2, \dots$. In round $d$ it visits the
diagonals $k = -d, -d+2, \dots, d$. For each it takes the starting point
\begin{equation}\label{eq:start}
x_{\mathrm{start}}(d,k) = \max\bigl\{\, x(d-1,k-1) + 1,\; x(d-1,k+1) \,\bigr\}
\end{equation}
and then slides along matches as far as possible, which gives $x(d,k)$. It stops
when a slide reaches $(N_a, N_b)$.

\begin{lemma}[Classical~\cite{myers1986}]\label{lem:myers}
$x(d,k)$ is the largest $x$ on diagonal $k$ that is reachable from the origin
with exactly $d$ non-diagonal steps.
\end{lemma}

\begin{proof}[Proof sketch]
Any path to diagonal $k$ with $d$ non-diagonal steps makes its last such step
from diagonal $k-1$ or $k+1$ after $d-1$ of them. Replacing its prefix by the
furthest-reaching path on that diagonal does not move the landing point of the
step backwards, and sliding along matches afterwards only moves it forwards.
Induct on $d$.
\end{proof}

It follows that the algorithm stops in round $D$.

\begin{proposition}[Updates]\label{prop:updates}
The number $U$ of (round, diagonal) updates satisfies
\[
\frac{D(D+1)}{2} + 1 \;\le\; U \;\le\; \frac{(D+1)(D+2)}{2}.
\]
\end{proposition}

\begin{proof}
Round $d$ has $d + 1$ diagonals. Rounds $0, \dots, D-1$ run in full, which is
$\sum_{d=0}^{D-1}(d+1) = D(D+1)/2$ updates. Round $D$ makes at least $1$ and at
most $D + 1$.
\end{proof}

\begin{proposition}[Comparisons]\label{prop:compares}
The number $C$ of base comparisons satisfies
\[
C \;\le\; U + (2D+1)\min(N_a, N_b).
\]
\end{proposition}

\begin{proof}
Each slide ends at its first mismatch or at the edge of the grid, so there is at
most one failed comparison per update, and at most $U$ in total.

A successful comparison advances the slide by one point on its diagonal. We show
that no point is passed twice. Let $(x,y)$ be the point where the slide on
diagonal $k$ stopped in round $d - 2$. A vertical step from it lands on diagonal
$k - 1$ with the same $x$, so $x(d-1,k-1) \ge x$. A horizontal step lands on
diagonal $k + 1$ at $x + 1$, so $x(d-1,k+1) \ge x + 1$. If one of these steps
would leave the grid, the other still gives the bound, or the starting point
falls outside the grid and the slide makes no comparison. By \eqref{eq:start},
$x_{\mathrm{start}}(d,k) \ge x + 1$: each slide begins beyond the point where
the previous slide on the same diagonal stopped.

At most $2D + 1$ diagonals are ever visited, and a diagonal has at most
$\min(N_a, N_b)$ points.
\end{proof}

\begin{corollary}\label{cor:wfcost}
The worst-case work of the wavefront algorithm is
$O\bigl(D\min(N_a,N_b) + D^2\bigr)$.
\end{corollary}

Both bounds were checked on $3{,}000$ random pairs with no violation. For highly
similar sequences the observed work is close to $N + D^2$, because almost all
matches lie on one diagonal. That is an empirical observation, with an
expected-case analysis in the literature. It is not a worst-case bound.

\section{The segment-vote band}

Split the read into consecutive segments of length at most $\Sseg$. Segment $t$
votes for a diagonal $d_t$. A row in segment $t$ computes the diagonals from
$\min(d_{t-1},d_t,d_{t+1}) - w$ to $\max(d_{t-1},d_t,d_{t+1}) + w$, where a
missing neighbour is omitted. Let $G = \sum_t \abs{d_{t+1} - d_t}$.

\begin{proposition}\label{prop:segment}
The number of cells computed is at most $mW + 2\Sseg G$.
\end{proposition}

\begin{proof}
A row in segment $t$ computes at most $R_t + W$ cells, where $R_t$ is the range
of $\{d_{t-1}, d_t, d_{t+1}\}$. The range of three numbers is at most the sum of
the two consecutive absolute differences, so
$R_t \le \abs{d_t - d_{t-1}} + \abs{d_{t+1} - d_t}$. Summing over the at most
$\Sseg$ rows of each segment and over all segments,
\[
\text{cells} \;\le\; mW + \Sseg \sum_t R_t \;\le\; mW + 2\Sseg G ,
\]
since each difference $\abs{d_{t+1} - d_t}$ appears in at most two of the $R_t$.
\end{proof}

With no indel, $G = 0$ and the count is at most $mW$. With one indel of length
$g$, $G = g$. In the worst case the band covers the whole matrix.

\begin{remark}
This bounds the work only. The segment-vote band missed $7$--$13\%$ of
$5$--$12$~bp indels in simulation, all within one segment length of the end of
the overlap, and there is no correctness theorem for it.
\end{remark}

\section{The gate}

\begin{proposition}\label{prop:gate}
If a read passes the gate with probability $\rho$, the vote costs
$C_{\mathrm{vote}}$ and the window costs $C_{\mathrm{SW}}$, then the expected
cost per read is
\[
\mathbb{E}[T] = C_{\mathrm{vote}} + (1 - \rho)\,C_{\mathrm{SW}},
\qquad
C_{\mathrm{vote}} \le \mathbb{E}[T] \le C_{\mathrm{vote}} + C_{\mathrm{SW}} .
\]
\end{proposition}

\begin{proof}
Every read pays for the vote. The window runs with probability $1 - \rho$. The
claim follows by linearity of expectation.
\end{proof}

The value of $\rho$ is empirical. It depends on the threshold, which has not
been chosen, and on the reads.

\section{What is not proved}\label{sec:open}

\begin{enumerate}
\item \textbf{$S_{\W} = S$ in general.} Corollary~\ref{cor:contain} gives
      $\score(T) \le S_{\W} \le S$. Full Smith--Waterman may find a
      higher-scoring alignment that leaves the window, and nothing here excludes
      that. In simulation it happened in about $1$ read in $16{,}500$.
\item \textbf{The vote for reads outside the model.} Theorem~\ref{thm:vote}
      covers one indel of length at most $w$ together with substitutions. Reads
      with two or more indels, chimeric reads, and reads from elsewhere in the
      genome are not covered.
\item \textbf{The vote with $\lambda = 0$.} The margin is negative for some
      templates, so there is no guarantee.
\item \textbf{The margins on the final windows.} Section~\ref{sec:margins} used
      provisional window coordinates and forward-strand reads. The numbers must
      be recomputed once the coordinates and the strand handling are fixed.
\item \textbf{The gate.} There is no theorem about which reads may skip
      Smith--Waterman. In simulation, $3.6\%$ of indel reads had an ungapped
      score equal to the full score because the indel was within $14$~bp of the
      end of the overlap.
\item \textbf{The segment-vote band and the adaptive band.} Only the work bound
      of Proposition~\ref{prop:segment} is proved. The adaptive band was worse
      than a fixed band in simulation.
\item \textbf{Typical-case cost of the wavefront method.}
      Propositions~\ref{prop:updates} and~\ref{prop:compares} are worst-case
      bounds only.
\item \textbf{Speed and power of real hardware.} Table~\ref{tab:hw} counts cells
      and cycles. Clock frequency, routing and power have not been measured.
\end{enumerate}

\section*{Scripts}

\begin{center}
\begin{tabular}{@{}>{\raggedright\arraybackslash}p{0.47\textwidth}>{\raggedright\arraybackslash}p{0.47\textwidth}@{}}
\toprule
Result & Script \\
\midrule
Margins (Tables~\ref{tab:margins}, \ref{tab:shares}) & \code{reference/vote_margin.py} \\
Difference ranges (Corollary~\ref{cor:bits}) & \code{reference/ideas_study.py} \newline \code{reference/window_study.py} \\
Window and vote on simulated reads & \code{reference/band_study.py} \\
Wavefront counts (Propositions~\ref{prop:updates}, \ref{prop:compares}) & \code{reference/ideas_study.py} \\
\bottomrule
\end{tabular}
\end{center}

\begin{thebibliography}{9}

\bibitem{smith1981}
T.~F. Smith and M.~S. Waterman.
\newblock Identification of common molecular subsequences.
\newblock \emph{Journal of Molecular Biology}, 147(1):195--197, 1981.

\bibitem{myers1986}
E.~W. Myers.
\newblock An {$O(ND)$} difference algorithm and its variations.
\newblock \emph{Algorithmica}, 1:251--266, 1986.

\bibitem{fischer1974}
M.~J. Fischer and M.~S. Paterson.
\newblock String-matching and other products.
\newblock In \emph{Complexity of Computation}, SIAM--AMS Proceedings, volume~7,
  pages 113--125, 1974.

\end{thebibliography}

\end{document}
